\section{Matching states and structural operations}
\label{sec:matching-state}

We describe a perfect-matching procedure on a loopless graph of maximum
degree three.  This section develops its combinatorial state, exact
ledger transformations, and structural searches.  The local implementation
of projection and the phase analysis follow in
Section~\ref{sec:projection-analysis}.  The dynamic optimization interface
used here is Lemma~\ref{lem:priced-steps}, proved later.

The graph has \(N\le S\) vertices, where \(S=2^s\) is a master size bound.
Disconnected graphs are allowed; parallel edge copies retain separate
identifiers, and degree counts incidences.  A successful answer is an
explicit perfect matching, checked before acceptance.  A negative answer
may arise from a capped randomized search, whose error is bounded later.
A deterministic prerequisite failure instead suspends this parameter
branch without an answer.  Perfect matchability will prove progress,
never a running-time bound.

\subsection{The phase state and its ledgers}
\label{match:parameters}

The empty graph is settled directly, and odd \(N\) gives a negative
answer.  Otherwise start with the empty matching and work in phases.
A phase starts with an actual matching \(M\), whose positive deficit
is \(K=N/2-|M|\).  At deficit zero the procedure is finished.
Within a phase we represent the changing matching by odd blocks and
single vertices assigned to those blocks.

A \emph{bin} is a physical odd vertex set with a constructive
factor-critical witness: after prescribing any one vertex as its
exposure, the others can be matched internally.  Its weight is
\(w_d=(|d|+1)/2\).  The bins form a disjoint family \(\mathcal D\).
The \emph{rows} \(\mathcal A\) are physical singleton vertices outside
them.  An explicit injection
\[
                         \mu:\mathcal A\longrightarrow\mathcal D
\]
assigns each row along a specified physical edge to its bin.
The bipartite incidence graph \(\mathcal B\) contains \emph{every}
physical row--bin edge, with different incidences retained separately.
Each row has degree at most three.

An assigned bin exposes the physical endpoint of its assigned edge
and is matched there to its row.  A free, or unassigned, bin exposes
one unmatched vertex.  The remaining unblocked vertices form a set
\(U\) of explicitly matched pairs.  A phase may also set aside
individual unmatched vertices as \emph{blockers}.  This state specifies
an actual matching, although a changed internal exposure need not be
propagated through a bin witness until materialization.  Initially the
bins are the unmatched singleton vertices, \(U\) contains all matched
pairs, and there are no rows or blockers.

Constructive odd blocks, with an exposed vertex chosen during expansion,
follow the blossom contraction-and-expansion tradition of
Edmonds~\cite[Sections~4.10--4.14 and 7.1--7.2]{EdmondsPaths}.
Here each block carries the witness needed to expand the represented
matching; the phase analysis below proves the required progress bound.

For \(X\subseteq\mathcal D\), define its captive rows and deficiency by
\begin{equation}
 A(X)=\{r\in\mathcal A:N_{\mathcal B}(r)\subseteq X\},
 \qquad h(X)=|X|-|A(X)|.
 \label{match:deficiency}
\end{equation}
Each row has its assigned bin as a neighbor.  The injection therefore
gives \(h(X)\ge0\), even at intermediate states.
Write \(v(X)=\sum_{d\in X}v_d\) for a bin function \(v\).

Fix integers \(a\ge1000\) and \(0\le\sigma\le a+1\), and set
\(C=a+1\).  The phase seeks
\begin{equation}
 g_*=\max\left\{1,\left\lfloor\frac{K}{100C}\right\rfloor\right\},
 \qquad L=\left\lceil\frac{10000CN}{K}\right\rceil
 \label{match:phase-parameters}
\end{equation}
gains, each increasing the matching size by one.  The parameter \(L\)
will be the amount allocated by each row in the ledgers below.
The choices balance the gain target against the work permitted per
gain: \(Lg_*=O_a(N)\), including when \(K<100C\), since then
\(g_*=1\) and \(L\le10000CN+1\).  At a phase boundary we materialize
the matching and initialize a new state.

For the reset periods, pricing steps and exact arithmetic, put
\begin{equation}
 B=2^{\lceil s/a\rceil},\qquad q=B^{-100},\qquad
 \delta=S^{-1000},\qquad \varepsilon=S^{-3000}.
 \label{match:small-parameters}
\end{equation}
The top ledger uses the grid \(\varepsilon\mathbb Z\); the other
ledgers use integers.  Their common purpose is to retain enough
unallocated capacity at each bin to rearrange the row assignment.
Debt records a temporary shortfall from the specified residual floor.

\begin{definition}[Ledger invariant]
\label{match:invariant}
There are ledgers \(j=0,\ldots,a\), with \(c_j=C-j\).
Ledger \(j\) has nonnegative incidence allocations \(x^j_e\), with
\[
 \sum_{e\text{ at row }r}x^j_e=L\qquad(r\in\mathcal A).
\]
Its column residual and nonnegative debt are
\[
 y^j_d=L-\operatorname{load}^j_d,\qquad
 \operatorname{load}^j_d=\sum_{e\text{ at bin }d}x^j_e,\qquad
 p^j_d\ge0,
\]
and they satisfy
\begin{equation}
                         y^j_d+p^j_d\ge c_jw_d.
 \label{match:ledger-floor}
\end{equation}
At the top level,
\begin{equation}
 x_e=x^0_e>0,\qquad s_d=y^0_d+p^0_d-Cw_d>0
 \label{match:top-slack}
\end{equation}
are \(\varepsilon\)-grid coordinates.  The other ledgers are integral.
At each settled main-loop state, the bottom debt \(p^a\) is zero.
\end{definition}

Initially allocations are empty and debts vanish.  Singleton residuals
are \(L\), and \(L>C\), so the top slacks are positive.

\begin{lemma}[Bottom surplus]
\label{match:bottom-surplus}
At each settled state and every nonempty \(X\subseteq\mathcal D\),
\begin{equation}
                       Lh(X)\ge y^a(X)\ge w(X)>0.
 \label{match:bottom-certificate}
\end{equation}
Thus \(h(X)\ge1\), every bin has weight at most \(L\), and a tight
set \(h(X)=1\) has total weight at most \(L\).
\end{lemma}
\begin{proof}
Every captive row allocates all \(L\) units inside \(X\), so
\(\operatorname{load}^a(X)\ge L|A(X)|\) and \(y^a(X)\le Lh(X)\).
Zero bottom debt and \(c_a=1\) give \(y^a(X)\ge w(X)\).
Weights are positive integers, proving the assertions.
\end{proof}

The ledgers have three roles.  At settled states the debt-free bottom ledger supplies the
strict Hall surplus just proved; the top ledger supplies a smooth progress
objective; and the intermediate integral ledgers make repairs local and
amortizable.  Adjacent floors differ by one.  In a failed projection,
this gap provides the margin that pays for the peeled
weight in each preserved coarser ledger, as proved in
Lemma~\ref{match:exchange-projection}.

\subsection{The top objective and the pricing primitive}

The top ledger chooses its next allocation by increasing
\begin{equation}
             \Phi=\sum_{e\in E(\mathcal B)}\log x_e+
                         \sum_{d\in\mathcal D}\log s_d.
 \label{match:top-objective}
\end{equation}
Orient allocation arcs from rows into bins, and add a slack arc
\(s_d\) from one common auxiliary vertex into each bin.
For fixed \(p^0,w\), circulation increments preserve exactly the
row sums and the column equations
\[
                 \sum_{e\text{ at }d}x_e+s_d=L+p^0_d-Cw_d.
\]
Conversely, respecting these equations gives a circulation:
the common slack-source equation follows by summing them and the
row equations.

Here is the part of Lemma~\ref{lem:priced-steps} needed to read the
matching algorithm.  Let \(f\) denote the vector of its positive
allocation and slack coordinates; here the gradient and derivative-sum
lengths are \(g_e=\lambda_e=1/f_e\).  A successful pricing call makes
an exact \(\varepsilon\)-grid circulation increment \(h\) with
\(\sum_e |h_e|/f_e\le q/100\), preserves positivity, and increases
\(\Phi\) by \(\Omega(q^2)\).  Its alternative outcome
\(\mathrm{LOW}\) is impossible whenever a circulation direction
\(d\ne0\) satisfies
\[
                  \sum_e\frac{d_e}{f_e}
                     \ge\frac12\sum_e\frac{|d_e|}{f_e},
\]
unless the primitive reports a stochastic failure or suspends.
The sums here range over both kinds of arc in the top graph.

The primitive stores the exact values implicitly and supplies cached
registers \(\widehat f_e\).  After settling its notifications they obey
\( |f_e-\widehat f_e|\le(q/100)\widehat f_e\).
It reports refreshed coordinates in a list, so the client can repair
small arguments and update expansion eligibility without scanning all
coordinates after a step.  The number of these notifications is charged
to successful pricing steps by Lemma~\ref{lem:priced-steps}.
Explicit point reads and settings remain available.  Section~\ref{sec:engine}
proves this interface and its time and reliability bounds; below we verify
the positivity and representation conditions required of this client.

Check the loose prerequisites
\begin{equation}
 S\ge2^{100},\qquad 100000C\le S^{1/20},\qquad
 B\le S^{1/20},\qquad q\ge S^{-1}.
 \label{match:initial-guards}
\end{equation}
Cap main-loop actions, including structural operations and pricing
attempts, at \(S^{20}\) per phase.  Separately cap pricing iterations in
each local threshold test of Section~\ref{sec:projection-analysis}
at \(S^{20}\).  Each engine instance has its stated \(S^{100}\)
external-operation guard, together with incrementally checked numerical
and representation bounds.  Failed prerequisites or deterministic caps
suspend the branch.  Exhaustion of a prescribed randomized search
instead returns a negative perfect-matching answer immediately.
The guarded-prefix analysis will show that the deterministic checks
interrupt no path for sufficiently large \(S\), at each fixed
sufficiently large choice of \(a,\sigma\).

\subsection{The main actions and exact ledger edits}

The actions below distinguish progress in matching size from changes
to its representation.  A gain increases the size by one; contraction
and expansion preserve it.  A projection may also peel bins, moving
their free exposures to the blocker list without changing matching
size.  Whenever the phase continues, the scheduled projection restores
the debt-free bottom ledger before the next settled state.

At a settled state, take the first applicable action in this order.
\begin{enumerate}[label=\arabic*.]
\item If a physical edge joins distinct bins \(i,j\), use it.
If the injection can be rearranged to leave both free, match their
prescribed exposed endpoints along this edge and move their vertices
into \(U\).  This is a gain.  The \(g_*\)-th gain ends the phase
immediately, without further ledger repair.
Otherwise contract the unique minimal tight \(X\supseteq\{i,j\}\),
with \(h(X)=1\), together with its captive rows to one bin.
The witnesses and efficient randomized implementation are proved below.
\item If bin \(d\) has an edge to \(v\in U\) and cached top slack
\(\widehat s_d\ge4C\), expand along it.  Make \(v\) a row and its
current \(U\)-mate a singleton bin, with their old match as assigned edge.
\item If neither action applies, request a top pricing step from
Lemma~\ref{lem:priced-steps}.  LOW or explicit sampling failure returns
a negative answer; suspension suspends the parameter branch.
\end{enumerate}
After each nonterminal structural operation, perform the scheduled
projection below before returning to a settled state.

In a contraction, \(|A(X)|=|X|-1\), so the new block has physical
size \(2w(X)-1\) and weight \(w(X)\).  Delete captive rows and preserve
every other assignment.  At most one exterior assignment entered
\(X\), because captive rows already assigned \(|X|-1\) of its bins.
Retain allocations on surviving incidences and aggregate \(p,y\),
and top \(s\), by summing over \(X\).  This is exact: removing captive
rows deletes \(L(|X|-1)\) load, while \(|X|\) residual columns become
one.  New incidences, including edges from a surviving row to a
swallowed row, start at zero in integral ledgers; the top uses the
positivity modifications below.

On expansion, ledger \(j\ge1\) allocates the new row's whole \(L\)
to its assigned child and gives that child debt \(c_j\).
Its other new incidences start at zero.  At the top, place \(C+1\)
units on the trigger incidence into the parent and \(L-C-1\) on
the assigned child incidence, with no initial child debt.
The child has top slack one.  Lemma~\ref{match:top-maintenance}
ensures actual parent slack at least \(3C\), leaving it positive
after paying \(C+1\).  Other new top incidences receive the small
positive padding specified there.

The following deletion transformation applies to gains and projection
peels.  Given bins \(R\), remove all rows assigned to \(R\).
For a peel, materialize each assigned bin with its row as matched
pairs in \(U\).  For each free bin, place its exposed vertex on the
blocker list and its other vertices into matched pairs in \(U\).
At a gain the two newly freed bins instead use their joining edge,
so no blocker is introduced.

\begin{lemma}[Deletion and net debt]
\label{match:deletion}
In a preserved ledger of floor \(cw\), redirect each allocation of
a surviving row into \(R\) to its current assigned edge, adding an
equal debt at the recipient bin.  This preserves the exterior
ledger invariant.  If \(O\) is redirected mass, \(I\) is removed
rows' allocation to the exterior, and \(f\) is the number of free
bins removed, then
\begin{equation}
 \begin{split}
 O&=I+Lf-y(R),\\
 I&\le L(h(R)-f),\\
 \Delta\!\left(\sum_d p_d\right)=O-p(R)&\le Lh(R)-cw(R).
 \end{split}
 \label{match:deletion-debt}
\end{equation}
The last bound is at most \(2L\) at a gain.
A gain decreases the represented unmatched count by two;
contractions, expansions, and peels leave it unchanged.
\end{lemma}
\begin{proof}
A surviving row's assigned bin is outside \(R\).
The added allocation and debt cancel in \(y+p-cw\);
removing rows only increases exterior slack.

Exactly \(|R|-f\) rows are removed.  Their allocation into \(R\)
is \(L(|R|-f)-I\), giving
\[
 y(R)=L|R|-\bigl(L(|R|-f)-I+O\bigr)=Lf+I-O.
\]
Every captive row is removed, as its assigned bin lies in \(R\).
Only the other \(h(R)-f\) removed rows contribute to \(I\).
Deleting debt \(p(R)\) and adding \(O\) gives the exact net debt
change, and \eqref{match:ledger-floor} gives the inequality.
At a gain, \(R\) has two bins and \(h(R)\le2\).

There is one unmatched vertex per free bin or blocker.
Expansion adds one bin and one assigned row.  Contraction replaces
\(|X|\) bins and \(|X|-1\) assigned captive rows by one bin,
preserving any exterior assignment.  A peel transfers free bins'
unmatched vertices to blockers.  A gain matches two free exposures.
This proves the unmatched-count assertions.
\end{proof}

\subsection{Maintaining the top barrier}

Structural edits change the graph and its column equations, while pricing
preserves those equations.  We now specify the local positivity repairs
that make both operations compatible with the pricing interface.

\begin{lemma}[Top maintenance]
\label{match:top-maintenance}
The top ledger can be maintained on the \(\varepsilon\)-grid with
logarithmic arguments above \(\delta\) before pricing and above
\(\delta/10\) at refreshes.  Apart from structural work, repairs
cost constantly many explicit point operations per dirty coordinate.
Between structural operations a small-allocation repair decreases
\(\Phi\) by at most \(O(\delta)\), and slack repairs never decrease it.
Positivity repairs add less than \(S^{-800}\) debt on every guarded
phase prefix.  Moreover,
\begin{equation}
 \begin{split}
 \widehat s_d\ge4C&\ \Longrightarrow\ s_d\ge3C,\\
 \text{no eligible expansion}&\ \Longrightarrow\
 s_d\le5C\quad\text{for every bin with a \(U\)-contact}.
 \end{split}
 \label{match:contact-slacks}
\end{equation}
\end{lemma}
\begin{proof}
A new positive allocation receives \(4\delta\) from a large donor
in its row.  Reading its at most three entries finds such a donor,
since their sum is \(L\).  Add \(4\delta\) debt at the destination
bin.  Its increased load is canceled by the debt; the donor bin's
slack increases by \(4\delta\).  If both bins coincide, their load
is unchanged and their slack still increases by that amount.
The two specified large entries of a newly expanded row may be
installed first and used as donors for its remaining incidences.

Use the same transfer whenever an existing allocation has cached
value at most \(3\delta\).  If a cached slack is at most \(3\delta\),
increase that slack and its debt by \(4\delta\).  Read touched values
exactly and re-anchor explicit settings.  Check touched slacks after
exact re-anchoring too, including values at the small-value threshold.

Between motions or structural changes, each row has at most three
entries to repair.  A donor initially at least \(L/3\) remains large
after these tiny transfers.  Both changed allocations are re-anchored
at nonsmall values.  Transfers never decrease a slack, and a slack
repair changes no other barrier coordinate.  Engine notifications
arising during settlement are queued and served, but price and
forest maintenance itself causes no primal motion.  Without a new
motion or structural change, an already repaired coordinate cannot
require another repair.  Thus there is no uncharged repair chain.

For a single positive logarithm the register bound gives
\[
                         |f-\widehat f|\le(q/100)\widehat f.
\]
After service, an unrepaired coordinate has cached value above
\(3\delta\), while a repaired one has been increased and re-anchored
by \(4\delta\).  This bound and the step-size bound leave the stated
margins before a step and at refreshes after it.
Contraction aggregates positive slacks; deletion increases surviving
slacks; expansion has the positive margins specified above.
New descriptors are inserted only after their positive final values
are computed.  All settings are exact \(\varepsilon\)-grid operations.

Only the large donor's logarithm can decrease during a transfer,
by \(O(\delta)\); recipient allocation and donor slack increase.
A slack repair increases its logarithm.  There are \(O(S)\)
coordinates at any state, with constantly many repairs per coordinate
or changed physical incidence between motions or structural batches.
The \(S^{20}\) phase-action cap and the no-motion observation give,
very crudely, at most \(S^{30}\) such fixes on a guarded prefix,
including the last completed edit batch.  Consequently their
\(4\delta\) debt additions total less than \(S^{-800}\).
This precision estimate does not permit scanning all coordinates
after each step: only touched entries and dirty notifications are
serviced.

Maintain bin-to-\(U\) contact lists and a cached-slack eligibility
index.  The relative register inequality, with \(q<1\), gives
\(s_d\ge3C\) when \(\widehat s_d\ge4C\); if no eligible contact
exists, its cached slack is below \(4C\) and its true slack below
\(5C\).
\end{proof}

An explicit edit batch reads affected coordinates exactly, retains
unchanged allocations on retargeted incidences, and computes positive
final settings before insertions or descriptor changes.  No pricing
occurs before its full row and column equations settle.  The identity
\(y^0_d=s_d-p^0_d+Cw_d\) supplies a top residual without summing
implicitly moving incidences.  Old descriptors and edge mappings
remain available through removal transactions.  Dirty-entry service
is completed before selecting the next main action.

Use nominal top-engine lifetime size \(R=SB\).  Lifetime size counts
vertex and edge births and endpoint insertions, including retargeting;
mere value settings, positivity fixes, and repricing on retained
topology do not create new incidences.  Lemma~\ref{lem:priced-steps}
charges
\[
 \bigl(R+B+\#\{\text{external pricing, point, or edge operations}\}\bigr)
                  S^{O(1/a)+o_a(1)}
\]
instructions under \(R+B\le S^{1+2/a}\), with polynomial labels
renumbered if needed.  The phase analysis verifies this nominal
size and the numerical bounds on sufficiently large guarded prefixes.
Engine notifications are paid by actual pricing motion, not by a
global scan.

\subsection{Reset schedule and the exact exchange problem}
\label{match:reset}

A phase has a tick counter initially zero.  A contraction or expansion
contributes one tick; a gain needing ledger repair contributes \(L\).
A projection peel contributes no ticks.  For \(j=0,\ldots,a\), put
\begin{equation}
                             b_j=B^{a-j}.
 \label{match:periods}
\end{equation}
After preliminary structural edits, choose the smallest
\(j\in\{1,\ldots,a\}\) whose period quotient increased during the
event.  Such a level exists since \(b_a=1\).  Construct a fresh
zero-debt \(j\)-ledger of floor \(c_jw\), and give every deeper
level a copy with zero debt.  Projection may peel one simultaneous
bin set.  The coarser levels \(0,\ldots,j-1\) are preserved through
it using Lemma~\ref{match:deletion}.

For \(j>1\), start from the just-current integral parent \(j-1\),
whose debt total is \(P=\sum_dp_d\).  For \(j=1\), instead use
\(L\) on each current assigned edge and zero elsewhere; no top
snapshot is needed.  Set \(c=c_j\) and define
\begin{equation}
 b(d)=(cw_d-y_d)_+,\qquad u(d)=(y_d-cw_d)_+,\qquad
 \beta=\sum_db(d).
 \label{match:exchange-demands}
\end{equation}
The exchange network has source-to-bin capacities \(b(d)\),
bin-to-sink capacities \(u(d)\), and, for every incidence \(e=(r,d)\),
\[
                d\longrightarrow r\quad[x_e],
       \qquad   r\longrightarrow d\quad[L+1],
\]
where brackets indicate integral upper capacities and all lower
capacities are zero.

\begin{lemma}[Exact exchange projection]
\label{match:exchange-projection}
An integral maximum flow in the exchange network produces either
a zero-debt ledger of floor \(cw\) on the whole state, or a nonempty
bin set \(R\) whose peel leaves such a ledger on the remaining state.
In the latter case,
\begin{equation}
                              Lh(R)<c_jw(R).
 \label{match:failed-floor}
\end{equation}
Every preserved coarser ledger loses at least \(w(R)\) in net debt,
and the number \(f\) of new blockers satisfies
\begin{equation}
                         f\le h(R)<\frac{Cw(R)}{L}.
 \label{match:peel-blockers}
\end{equation}
\end{lemma}
\begin{proof}
For an integral maximum flow \(F\), set
\[
                   x'_e=x_e-F_{d,r}(e)+F_{r,d}(e).
\]
Every row has total incoming exchange capacity \(L\), so its
outgoing exchange flow is at most \(L\).  Thus \(x'\ge0\), its
row sums are \(L\), and no row-to-bin arc of capacity \(L+1\)
is saturated.  If \(F_{\mathrm{src},d}\) and
\(F_{d,\mathrm{snk}}\) denote source and sink arc flows, then
\[
                 y'_d=y_d+F_{\mathrm{src},d}-F_{d,\mathrm{snk}}.
\]
At value \(\beta\), all deficits are filled while a spare loses
at most \(u(d)\), proving \(y'\ge cw\).

Otherwise let \(R\) be \emph{all} residual-reachable bins.
For an incidence \(e=(r,d)\), residual capacity from \(d\) to \(r\),
combining the original arc and the reverse of its companion, is
\[
                       x_e-F_{d,r}(e)+F_{r,d}(e)=x'_e.
\]
A positive allocation into a reachable bin makes its row reachable.
All that row's neighbors are reachable because its row-to-bin arcs
are unsaturated.  The rows allocating into \(R\) are therefore
exactly its captives, which allocate their full \(L\) there.  Hence
\(y'(R)=Lh(R)\).  Every reachable spare sink arc is saturated,
whereas some reachable source deficit remains unsaturated.
Summing the residual changes over \(R\) gives
\eqref{match:failed-floor}.

Every deficit outside \(R\) is saturated, as otherwise its source
arc would make it reachable; so \(y'\ge cw\) there.  A surviving
row has no allocation into \(R\), since such a row would be captive
and removed.  Removing the assigned rows of \(R\) only increases
exterior residuals.  Restricted \(x'\) is the requested ledger.

For a coarser level \(\ell<j\), one has \(c_\ell\ge c_j+1\).
Lemma~\ref{match:deletion} gives net debt change at most
\[
                         Lh(R)-c_\ell w(R)<-w(R).
\]
Thus at least \(w(R)\) is paid.  Every captive is among the
removed assigned rows, giving \(f\le h(R)\); now
\eqref{match:failed-floor} and \(c_j\le C\) give the blocker bound.
\end{proof}

At \(j=1\), construct this network globally and use
Lemma~\ref{lem:integral-flow}, with
Lemma~\ref{lem:capacity-scaling} if necessary.  Its size is \(O(S)\)
and its capacities are polynomial in \(S\), yielding
\(S^{1+o(1)}\) time including the arc-flow vector and residual cut.
For \(j>1\), the parent inequality \(y+p\ge(c_j+1)w\)
permits the local implementation of Lemma~\ref{lem:projection}.
It returns these exact allocation modifications and, when needed,
the entire residual-reachable peel set in
\[
                            (P+B)S^{O(1/a)+o_a(1)}
\]
time.  Its engine calls have their individual conditional failure
budgets.  A flagged failure ends this perfect-matching test negatively;
a deterministic suspension suspends the branch.

\subsection{Tight sets and factor-critical contraction}

\begin{lemma}[The tight-set alternative]
\label{match:tight-contraction}
Suppose \eqref{match:bottom-certificate} holds and an edge joins
distinct bins \(i,j\).  Either the injection can be rearranged to
leave both free, or a unique inclusion-minimal
\(X\supseteq\{i,j\}\) satisfies \(h(X)=1\).
In the latter case \(w(X)\le L\), and the physical block
\(\bigl(\bigcup_{d\in X}d\bigr)\cup A(X)\) is constructively factor-critical using the joining
edge, its row--bin incidences, and its child-bin witnesses.
\end{lemma}
\begin{proof}
An injection avoiding forbidden bins \(F\) exists precisely when
\[
                        |A(Z)|\le |Z\setminus F|
                   \qquad(Z\subseteq\mathcal D).
\]
Necessity is immediate.  For sufficiency, apply Hall's condition~\cite[Theorem~1]{Hall1935}
to any row subset: if \(Z\) is its full neighbor set, those rows
belong to \(A(Z)\), and \(Z\setminus F\) is large enough.
For \(F=\{i,j\}\), the criterion becomes
\begin{equation}
                          h(Z)\ge |Z\cap\{i,j\}|.
 \label{match:joint-hall}
\end{equation}
By the bottom surplus, failure means exactly a tight set containing
both bins.

Each neighborhood-inclusion indicator in \(|A(Z)|\) is
supermodular; thus \(h\) is submodular.  Two tight sets containing
\(i,j\) have nonempty intersection and union, both of deficiency
at least one.  Submodularity forces their intersection to be tight.
Intersecting all such sets gives the unique minimal \(X\).
Its weight bound follows from Lemma~\ref{match:bottom-surplus}.

To expose a physical vertex \(v\) in child bin \(d\in X\),
match all captive rows into \(X\setminus\{d\}\).
The Hall inequalities are \(h(Z)\ge1\) when \(d\in Z\)
and \(h(Z)\ge0\) otherwise, for \(Z\subseteq X\).
Expose \(v\) inside \(d\) and each assigned edge's physical contact
inside every other child.

To expose a captive row \(r\), match \(i,j\) along the joining edge,
and match the other captives into \(X\setminus\{i,j\}\).
The needed Hall inequalities are
\begin{equation}
 h(Z)+\boldsymbol1_{\{r\in A(Z)\}}\ge |Z\cap\{i,j\}|
                         \qquad(Z\subseteq X).
 \label{match:captive-exposure-hall}
\end{equation}
A proper subset containing both \(i,j\) has \(h(Z)\ge2\) by
minimality and integrality.  Other nonempty proper subsets use
\(h(Z)\ge1\), and the empty set is harmless.  For \(Z=X\),
the deficiency is one and \(r\) is captive.  Thus
\eqref{match:captive-exposure-hall} holds.  Each child exposes its
actual prescribed physical contact, so the child witnesses lift
these injections and the joining edge to the desired internal
matchings.
\end{proof}

\subsection{Bounded searches for a bin pair}

We implement the alternative without examining the physical interiors
of visited bins, first testing for a small tight set and then trying
a joint freeing.

\paragraph{A capacity-one witness.}
Use a directed reverse graph with split vertices
\(d_{\rm in}\to d_{\rm out}\) of capacity one for each bin.
The source has capacity-two arcs to \(i_{\rm in},j_{\rm in}\).
A free bin's output has a capacity-two arc to the sink.
An assigned bin's output has capacity-two arcs to the inputs
of all neighbors of its assigned row, retaining individual incidences.
Original outdegrees are bounded.

For the minimal tight \(X\), its \(|X|-1\) captive rows are assigned
to all but one bin \(t\) of \(X\).  The set comprising the source,
all inputs of \(X\), and all outputs of \(X\setminus\{t\}\)
has at most \(2|X|\) vertices and outgoing capacity one: only
\(t\)'s split arc leaves it.  The original sink is reachable.
Otherwise all reachable bins are assigned, and their assigned rows
have every neighbor among those bins, contradicting positive
deficiency for that nonempty set.

Use doubling vertex budgets \(v_0=1,2,4,\ldots\), through the first
budget at least \(8L\).  A fresh trial explores the original graph
from the source until reaching the sink or visiting \(v_0\) vertices.
Send one unit along a predecessor path, ending at the sink if reached,
and otherwise at a uniformly sampled visited vertex.
Then explore residual reachability with the same budget.

When \(v_0\ge4|X|\), either the endpoint is the sink or exactly
\(v_0\) vertices were visited.  At most \(2|X|\) of the latter
lie in the small cut, so the endpoint is outside with probability
at least one half.  A simple predecessor path to the exterior
uses its sole outgoing arc once and cannot return: it would then
have to use that same outgoing arc again to reach its endpoint.
It therefore saturates the outgoing capacity and has no positive
incoming flow across the cut.  The cut has no outgoing residual arc.
Since \(v_0>2|X|\), residual search then exhausts without the sink.

Only a \emph{fully exhausted, sink-free} residual source side is
accepted.  Its original outgoing arcs are saturated and its incoming
flow is zero, so its original outgoing capacity equals net path flow,
at most one.  That capacity is positive since the original sink was
reachable.  It is therefore exactly one.  No capacity-two arc crosses
the source side: the unique crossing is a split arc, both specified
inputs are inside, and every included output has all of its assigned
row's neighbors inside.

An included output has its own input inside too.  It is reached either
through its split arc, or through the residual reverse of an original
outgoing arc.  In the second case that outgoing arc carried the
predecessor-path flow, so its split arc also carried flow and has
a residual reverse to its input.  Thus the included inputs form a
bin set \(X'\) in which all except the crossing-split bin have distinct
captive assigned rows.  Consequently \(h(X')\le1\), and the bottom
surplus makes \(X'\) a tight witness containing \(i,j\).
Truncated searches and searches reaching the sink certify nothing.

\paragraph{Refining a witness.}
List the captives of \(X'\) through its assigned bins, checking each
row's bounded neighbor list.  On these rows and bins, use source-to-row
capacities one, row-to-bin capacities \(H_\infty\), bin-to-sink
capacities one, and additional source arcs of capacity \(H_\infty\)
to \(i,j\), where \(H_\infty=2|X'|+3\) exceeds the finite cut values
of interest.  Given a bin source side \(Z\supseteq\{i,j\}\), the
optimal row source side is \(A(Z)\), with cut capacity
\[
                             |A(X')|+h(Z).
\]
The minimum is \(|A(X')|+1\).  Start with the captive assigned edges,
a flow of value \(|A(X')|\).  One augmentation reaches the maximum,
and residual reachability gives its smallest source-side minimum cut.
Its bins are the unique minimal tight \(X\).  This costs linear
work in \(|X'|\), up to dictionaries.

Repeat fresh trials at each budget as specified below.
After refinement, if the discovery budget exceeds \(16|X|+16\),
abort the perfect-matching test negatively.  On a tight pair this
requires failure of every trial at the first budget at least
\(4|X|\).  On every path, however, a continued contraction costs
only \(O(|X|\polylog S)\), with fixed-parameter reliability factors;
a final abort may cost \(O(L\polylog S)\).
The one pushed path creates only linearly many residual reversals,
so each bounded exploration has linear budget cost up to dictionaries.

\paragraph{Freeing two bins.}
If no tight witness is found, try to free the bins jointly.
At an assigned bin, a reverse walk follows an incidence of its
assigned row with probability \(x^a_e/L\); free bins absorb.
For the transition submatrix \(P_m\) on assigned bins, each
destination column \(d\) has sum at most
\[
                  \operatorname{load}^a_d/L\le1-w_d/L.
\]
This remains true after an injection rearrangement: each assigned
row still supplies exactly one transition row.  Since \(w_d\ge1\),
the probability of surviving \(\ell\) steps from any start is at most
\begin{equation}
                              N\exp(-\ell/L).
 \label{match:walk-survival}
\end{equation}
Indeed the matrix column norm is at most \(1-1/L\).
Each entry of \(P_m^\ell\) is at most its column sum, and summing
a starting row over at most \(N\) columns proves the bound.

First free \(i\) by a successful walk to a free bin: erase loops
while keeping the assignment fixed, and then shift each encountered
row to the next bin of the resulting path.  If \(i\) was already
free, no change is needed.  With the new injection, joint freeability
guarantees a path from \(j\) to another free bin, avoiding \(i\),
in the \emph{full-neighbor} reverse graph.  Otherwise the exhausted
set \(T\) of reached assigned bins has all its assigned rows captive
in \(T\cup\{i\}\).  This set contains \(i,j\) and has deficiency at
most one, contrary to \eqref{match:joint-hall}.

Explore until another free bin is reached, or until reaching \(8L\)
distinct assigned bins.  In the first case use its predecessor path.
In the second, sample a reached bin uniformly, retain its prefix
path from \(j\), and try a bottom-ledger walk from there.
Let \(H_d\) be the probability of eventually hitting \(i\) from an
assigned bin, and \(p_i\) the one-step column into \(i\).
For the assignment fixed during this attempt,
\begin{equation}
 H=P_mH+p_i,\qquad
 w^{\mathsf T}H\le L\boldsymbol1^{\mathsf T}(I-P_m)H
       =L\boldsymbol1^{\mathsf T}p_i\le L.
 \label{match:avoid-first-free}
\end{equation}
The first inequality is the destination-column bound with \(H\ge0\);
the last uses \(\operatorname{load}^a_i/L\le1\).
As weights are at least one, the average hit probability over
the \(8L\) reached bins is at most \(1/8\).
Together with \eqref{match:walk-survival}, this gives a constant
chance of reaching another free bin within the walk cap while
avoiding \(i\).  Concatenate the fixed prefix with that walk,
erase loops, and shift assignments.  This frees \(j\) while
leaving \(i\) free.

Loop erasure uses a stack indexed by bin.  Each retained outgoing
transition follows the row assigned at its origin, since the
assignment stays fixed throughout path construction.  Hence a path
to a free bin certifies the reassignment.  An already-free target
needs no rearrangement.  Exhaustion or failure after the capped
repetitions gives a negative answer; successful joint freeing
always certifies a genuine gain.

\begin{lemma}[A bounded randomized pair action]
\label{match:pair-action}
At a settled state, the preceding searches process a queried
distinct-bin edge as follows.  A returned contraction is the minimal
tight set with a valid factor-critical witness; a returned gain has
a valid joint-freeing witness.  Conditional on the valid current
state, the probability of a premature negative answer, including
a late-witness abort, is at most \(2^{-\sigma}S^{-60}\).
Continued contractions cost \(O(|X|\polylog S)\); a gain attempt
or final abort costs \(O(L\polylog S)\), with fixed-parameter
reliability factors.  These work bounds hold on every path.
\end{lemma}
\begin{proof}
Use walk length and repetition count
\[
                    \ell=16L(1+s),\qquad
                    T_{\rm rep}=1000(\sigma+s+1).
\]
Use \(T_{\rm rep}\) fresh tight-set trials at each budget, that
many attempts at first freeing, and that many sampled second-freeing
attempts.  A tight-set trial succeeds with probability at least
one half once its budget is at least \(4|X|\), so missing the
early witness is exponentially unlikely in \(T_{\rm rep}\).
For a jointly freeable pair, the first walk has constant success
probability by \eqref{match:walk-survival}.  The second has probability
at most \(1/8\) of hitting \(i\), by
\eqref{match:avoid-first-free}, and its nonabsorption probability
at the length cap is bounded by \eqref{match:walk-survival}.
Their sum is strictly below one.  Unsuccessful trials do not change
the injection, so the same conditional estimates apply repeatedly.

A range-uniform draw generates enough independent fair bits to
cover the range and rejects out-of-range strings; each elementary
attempt succeeds with probability at least one half.
Cap its string attempts at \(T_{\rm rep}\), returning negative
on exhaustion.  An integral allocation draw uses a uniform element
of \(\{1,\ldots,L\}\) and the row's at most three cumulative intervals.
Under \eqref{match:initial-guards}, fewer than \(S^{10}\) such
draws occur in one pair query, including every budget and trial.
A union bound over constant-probability trial failures and capped
fair-bit rejection failures gives \(2^{-\sigma}S^{-60}\), with
ample numerical slack.  Equivalently, couple unlimited exact
rejection to this capped process until its first cap failure.

The witness arguments prove soundness.  Doubling budgets sum to
a constant multiple of the last, so the post-refinement budget test
gives \(O(|X|\polylog S)\) for a continued contraction even on an
unfavorable random path.  Other attempts use \(O(L)\) vertex caps,
\(O(L\log S)\) walk caps, and the fixed repetition counts.

Generate outgoing reverse arcs from bin/assignee data and the
assigned row's physical neighbor list, of degree at most three.
Store path-created residual reversals and local visited records
in dictionaries.  No query scans all physical vertices of a visited
bin.  Random bits, initialization, loop erasure, record operations,
and explicit path or witness output are included in the work.
\end{proof}

\subsection{Explicit matching materialization}

\begin{lemma}[Materializing a bin]
\label{match:materialization}
A current bin can be materialized with any prescribed physical
exposure in time linear in its physical size up to logarithmic
factors.  All current bins can be dissolved at an erasure or phase
boundary within the corresponding total physical-size bound.
\end{lemma}
\begin{proof}
At a contraction store the child-bin records, captive rows and
their physical incidences, joining edge, and captive assignment
at contraction time.  This gives a contraction tree whose leaves
include both child-bin vertices and swallowed row vertices.
The stored captive assignment covers all but one child.

If the desired exposure \(v\) is in child \(d\), forbid \(d\) and
remove its captive assignment if present.  At most one row becomes
unmatched.  Lemma~\ref{match:tight-contraction} guarantees an injection
of all captives avoiding \(d\), so at most one ordinary row
augmentation suffices.  If \(d\) was the unused child, none is
needed.  Expose \(v\) inside \(d\), and the physical assigned-edge
contact in every other child.  A preserved exterior assignment
merely prescribes \(v\) to be its contact endpoint.

If the desired exposure is swallowed row \(r\), delete \(r\) and
its assigned edge, reserve \(i,j\) for the stored joining edge,
and remove any remaining assignments into those children.
At most two remaining rows become unmatched.  Inequality
\eqref{match:captive-exposure-hall} guarantees a matching of the
remaining rows into \(X\setminus\{i,j\}\), so at most two
augmentations suffice.  The joining edge prescribes exposures
in \(i,j\); every other child gets its row-edge contact.

In either case each child receives exactly one prescribed exposure
and is recursively materialized once.  The searches scan only this
node's stored abstract incidence graph, of size linear in its arity
because rows have degree at most three, rather than its children's
physical interiors.

At the start of dissolving a root, one traversal assigns subtree
leaf intervals.  The interval containing a prescribed physical
vertex identifies its direct child, with logarithmic search overhead
if needed.  No contraction is unary, and each record has size linear
in its arity.  The sum of record sizes and leaves is therefore
linear in the dissolved root's physical size.  Summing the local
searches, endpoint lookups, and single child visits proves the bound.
Disjoint roots at an erasure or phase boundary can be handled in
the same total physical-size time.
\end{proof}

Later structural accounting charges any subsequent return of these
vertices to bins as a new stay supplied by initialization or expansion.
Thus each dissolved contraction record pays for one materialization,
rather than repeated descent through an unchanged nested witness.
